\documentclass[11pt,a4paper,oneside]{book}
\usepackage[T1]{fontenc}
\usepackage[utf8]{inputenc}
\usepackage[english]{babel}
\usepackage{lmodern}
\usepackage{geometry}
\usepackage{microtype}
\usepackage{amsmath}
\usepackage[table]{xcolor}
\usepackage{hyperref}
\usepackage{bookmark}
\usepackage{longtable}
\usepackage{booktabs}
\usepackage{array}
\usepackage{enumitem}
\usepackage{listings}
\usepackage{fancyhdr}
\usepackage{graphicx}
\usepackage{float}
\usepackage{framed}
\usepackage{titlesec}
\geometry{margin=2.2cm}
\graphicspath{{images/}{docs/images/}}
\hypersetup{colorlinks=true,linkcolor=blue!55!black,urlcolor=blue!55!black,pdftitle={Optical equipment control panel manual}}
\setcounter{tocdepth}{2}
\setcounter{secnumdepth}{2}
\renewcommand{\arraystretch}{1.5}
\pagestyle{fancy}
\fancyhf{}
\usepackage{parskip}
\fancyhead[L]{Optical equipment control panel}
\fancyhead[R]{Operating and maintenance manual}
\fancyfoot[C]{\thepage}
\setlength{\headheight}{14pt}
\setlength{\tabcolsep}{4pt}
\setlength{\arrayrulewidth}{0.6pt}
\setlength{\doublerulesep}{0pt}
\definecolor{panelblue}{RGB}{21,62,117}
\definecolor{sectionblue}{RGB}{31,90,148}
\definecolor{tablehead}{RGB}{31,90,148}
\definecolor{codebg}{RGB}{246,248,250}
\definecolor{codeframe}{RGB}{210,215,220}
\titleformat{\chapter}[display]
  {\normalfont\huge\bfseries\color{panelblue}}
  {\chaptertitlename\ \thechapter}{12pt}{\Huge}
\titleformat{\section}
  {\normalfont\Large\bfseries\color{sectionblue}}
  {\thesection}{0.65em}{}
\titleformat{\subsection}
  {\normalfont\large\bfseries\color{sectionblue}}
  {\thesubsection}{0.65em}{}
\arrayrulecolor{panelblue!35}
\newcommand{\thead}[1]{\cellcolor{panelblue}\textcolor{white}{\textbf{#1}}}
\lstset{basicstyle=\ttfamily\footnotesize,backgroundcolor=\color{codebg},frame=single,rulecolor=\color{codeframe},breaklines=true,breakatwhitespace=true,columns=fullflexible,keepspaces=true,showstringspaces=false}
\DeclareRobustCommand{\code}[1]{\texttt{\detokenize{#1}}}
\Urlmuskip=0mu plus 1mu
\newcommand{\manualfigure}[2]{%
\begin{figure}[H]\centering
\IfFileExists{images/#1}{\includegraphics[width=0.94\textwidth,height=0.72\textheight,keepaspectratio]{images/#1}}{%
\fcolorbox{blue!45!black}{blue!4}{\parbox[c][5.5cm][c]{0.88\textwidth}{\centering\textbf{Figure missing:}\par\code{#1}\par\medskip #2}}}
\caption{#2}\label{fig:#1}\end{figure}}
\sloppy
\title{\Huge Optical equipment control panel\\[0.4cm]\Large Operating, administration, and maintenance manual}
\author{Hubert Makara \\ Michał Surówka}
\date{Date: 3 August 2026}
\begin{document}
\frontmatter
\maketitle
\tableofcontents
\mainmatter
\section{How to use this manual}

This manual provides step-by-step instructions for installing, configuring,
operating, and maintaining the amplifier panel. It explains the complete installation
process, including package installation, initial configuration, service startup,
and post-installation verification.

It also describes all statuses shown in the user interface, explains what data
and configuration files are stored on the device and where they are located,
identifies actions that modify the connected hardware or host system, and
provides troubleshooting procedures for common failures.

\begin{framed}
\textbf{Notice:} The software has not yet been validated against final hardware. Labels, serial responses, ranges, and state names may require adjustment.
\end{framed}
\chapter{System purpose and operating model}
The control panel is a local web application for monitoring and controlling one connected device. The selected \code{DEVICE_PROFILE} determines which protocol and interface are active:

\begin{longtable}{|>{\centering\arraybackslash}m{0.17\linewidth}|>{\centering\arraybackslash}m{0.31\linewidth}|>{\centering\arraybackslash}m{0.44\linewidth}|}
\hline
\thead{Profile} & \thead{Connected equipment} & \thead{What the control panel does} \\
\hline
\endhead
\code{amplifier} & Optical amplifier & Displays measurements, stores history, checks limits, and changes the gain setpoint \\
\hline
\code{fts-ls} & Frequency Transfer System laser station & Displays station and module states, stores history, and sends operator commands \\
\hline
\end{longtable}
Both profiles use the same authentication and roles, browser session, status bar, SQLite service, history controls, retention setting, CSV format, administration screens, diagnostics, file locations, backup procedure, and installation method. The serial protocol, displayed device fields, available commands, and validated alarm rules are profile-specific.

The main startup tasks, continuous acquisition loops, data processing, browser
refresh cycle, and user-action paths are shown in Figure~\ref{fig:11-program-operation-flow.png}.

\manualfigure{11-program-operation-flow.png}{Program operation flow, including device acquisition, browser refresh, warning processing, command handling, and network rollback loops.}
The browser never talks directly to the serial device. FastAPI owns the runtime state, a background worker owns the serial connection, SQLite stores history, and a separate restricted host agent performs NetworkManager changes.

Two systemd services are used:

\begin{description}[style=nextline]
\item[\code{amp-panel.service}] Runs the web interface, communicates with the connected device, stores history, and provides logging and SNMP functions. It runs as the unprivileged \code{amp-panel} user.
\item[\code{amp-panel-network-agent.service}] Applies restricted NetworkManager changes requested from the Administration screen. It runs separately because network configuration requires host privileges.
\end{description}
The web service is unprivileged and has no Linux capabilities. It receives only the supplementary groups required for serial devices and local logs.

\chapter{Login, roles, and navigation}
\section{Authentication model}
A username must pass two checks:

\begin{enumerate}
\item it must exist and be active in the local access list;
\item its password must be accepted by the configured RADIUS server.
\end{enumerate}
The control panel stores usernames, roles and active/inactive status. It does not store user passwords. The FTS-LS serial-console password is a device credential and is separate from the RADIUS login password.

Common login results are:

\begin{longtable}{|>{\centering\arraybackslash}m{0.20\linewidth}|>{\centering\arraybackslash}m{0.35\linewidth}|>{\centering\arraybackslash}m{0.35\linewidth}|}
\hline
\thead{Visible result} & \thead{Meaning} & \thead{Required response} \\
\hline
\endhead
Dashboard opens & Local access and RADIUS authentication succeeded & Continue according to the assigned role \\
\hline
Invalid username or password & Unknown/inactive local user or RADIUS rejection & Verify the exact username, local status, and RADIUS credentials \\
\hline
Authentication server unavailable & RADIUS could not be reached or its configuration is wrong & Check network, server address, UDP 1812, shared secret, and server status \\
\hline
Too many login attempts & Per-client rate limit was reached & Wait for the configured login window and investigate repeated failures \\
\hline
Login screen appears during use & Session expired, service restarted, or authentication cookie was lost & Sign in again \\
\hline
\end{longtable}
\section{Roles}
\begin{longtable}{|>{\centering\arraybackslash}m{0.23\linewidth}|>{\centering\arraybackslash}m{0.23\linewidth}|>{\centering\arraybackslash}m{0.23\linewidth}|>{\centering\arraybackslash}m{0.23\linewidth}|}
\hline
\thead{Capability} & \thead{Viewer} & \thead{Operator} & \thead{Administrator} \\
\hline
\endhead
Read live values & Yes & Yes & Yes \\
\hline
View history and statistics & Yes & Yes & Yes \\
\hline
Export measurement history & Yes & Yes & Yes \\
\hline
Change amplifier setpoint and warning thresholds & No & Yes & Yes \\
\hline
Acknowledge amplifier warnings & No & Yes & Yes \\
\hline
Operate normal FTS-LS controls & No & Yes & Yes \\
\hline
Manage users, network, NTP, services, and SNMP & No & No & Yes \\
\hline
Reboot, power reset, or factory reset the FTS-LS & No & No & Yes \\
\hline
\end{longtable}
Controls hidden by the browser are also protected by API role checks.

\section{Navigation groups}
\begin{itemize}
\item \textbf{Device} contains the live view and device settings.
\item \textbf{Monitor} contains Overview, Statistics, Warnings, and History.
\item \textbf{Admin} contains access control, SNMP, network, NTP, service diagnostics, and FTS-LS administrative actions.
\end{itemize}

The interface returns to the live view if a user opens a URL fragment for a page that their role cannot access.

\chapter{Status bar and common interface states}
The status bar remains visible above the selected page.

\section{Last update}
\code{Last update} is the timestamp of the most recent accepted device measurement or complete FTS-LS snapshot.

\begin{longtable}{|>{\centering\arraybackslash}m{0.28\linewidth}|>{\centering\arraybackslash}m{0.64\linewidth}|}
\hline
\thead{Display} & \thead{Meaning} \\
\hline
\endhead
Time value & At least one valid update was received \\
\hline
\code{--} & No valid device update has been accepted since process start \\
\hline
Old time while \code{CONNECTED} & The serial port may be open, but the device is not producing a complete supported message \\
\hline
\end{longtable}
An old timestamp is operationally important even if the connection indicator is green. Compare it with the system time and the expected polling interval.

\section{System time}
\code{System time} comes from the control-panel host through the live API and is
refreshed with each dashboard update. The browser formats the supplied timestamp
for display, but its own clock is not used as the time source. This value shows the
host clock; it is not independent proof that the clock is synchronized. Use
\textbf{Admin > Time diagnostics} for NTP information.

\section{Device status}
\begin{longtable}{|>{\centering\arraybackslash}m{0.20\linewidth}|>{\centering\arraybackslash}m{0.35\linewidth}|>{\centering\arraybackslash}m{0.35\linewidth}|}
\hline
\thead{State} & \thead{Exact meaning} & \thead{Typical causes} \\
\hline
\endhead
\code{CONNECTED} & The current serial session is open and the latest serial operation has not failed & Correct port and permissions; device is powered and responding \\
\hline
\code{DISCONNECTED} & No serial session is currently established & Cable/power issue, wrong port, port busy, permissions, wrong FTS-LS credentials, reconnect in progress \\
\hline
\code{API ERROR} & The browser could not obtain \code{/api/latest} & Web service restart, network interruption, expired session, server error \\
\hline
\end{longtable}
\code{CONNECTED} confirms the serial session. It doesn't guarantee correctness of displayed measurements. \code{API ERROR} describes communication between the browser and the control panel.

\section{Stored records}
This is the number of history records for the active profile. A value of zero can be valid immediately after installation. If live values change but the count never increases, inspect database diagnostics.

\section{Placeholder and loading states}
\begin{longtable}{|>{\centering\arraybackslash}m{0.28\linewidth}|>{\centering\arraybackslash}m{0.64\linewidth}|}
\hline
\thead{State} & \thead{Meaning} \\
\hline
\endhead
\code{--} & Value has not been reported, is not applicable, or the request has not completed \\
\hline
\code{Loading...} & A history, statistics, network, or diagnostic request is active \\
\hline
\code{Could not load data} & The request failed; existing live acquisition may still be operating \\
\hline
\code{Waiting for samples} & Too few records exist to estimate a write rate or retention period \\
\hline
\code{None} & The backend explicitly reports no current error/value \\
\hline
\end{longtable}
\chapter{Optical amplifier operation}
\section{Live measurements}
The panel shows values received from the amplifier together with values derived by
the application. It does not perform an independent optical measurement.

\begin{longtable}{|>{\centering\arraybackslash}m{0.20\linewidth}|>{\centering\arraybackslash}m{0.52\linewidth}|>{\centering\arraybackslash}m{0.18\linewidth}|}
\hline
\thead{Displayed field} & \thead{What it shows} & \thead{Unit} \\
\hline
\endhead
\code{PiA}, \code{PiB} & Optical power reported at input A and input B & dBm \\
\hline
\code{PoA}, \code{PoB} & Optical power reported at output A and output B & dBm \\
\hline
Actual gain & Current amplifier gain. The reported \code{gain_actual} value is used when available; otherwise the application calculates the mean gain of channels A and B from the four optical-power values & dB \\
\hline
Gain setpoint & Requested target gain. The device-reported value is used when available; otherwise the panel shows the last accepted setpoint saved by the application & dB \\
\hline
Gain delta & Gain setpoint minus actual gain. A positive value means the actual gain is below the target; a negative value means it is above the target & dB \\
\hline
Temperature & Temperature reported by the amplifier & \textdegree C \\
\hline
Sequence number & Measurement-frame counter supplied by the amplifier, when present & Unitless \\
\hline
\end{longtable}

When \code{gain_actual} is absent from a valid device message, the application uses:

\[
G_A = PoA - PiA, \qquad
G_B = PoB - PiB, \qquad
\text{Actual gain} = \frac{G_A + G_B}{2}.
\]

If \code{gain_delta} is absent, it is calculated as
\code{gain_set - gain_actual}. A missing value is displayed as \code{--}; the panel
does not replace a missing optical-power measurement with zero. The canonical API
and database field names remain \code{gain_actual}, \code{gain_set}, and
\code{gain_delta}. Raw firmware spellings are translated in
\code{app/protocols/amplifier.py}.

\section{Changing the gain setpoint}
\begin{enumerate}
\item Confirm \code{CONNECTED} and check the latest update time.
\item Record the current setpoint and actual gain.
\item Enter a value inside the safe range displayed by the interface.
\item Submit once.
\item Wait for a new device response and verify setpoint, actual gain, and delta.
\end{enumerate}
The server rejects non-numeric values, \code{NaN}, infinity, and values outside \code{GAIN_SET_MIN} to \code{GAIN_SET_MAX}. A request cannot succeed when the serial port is closed. The accepted setpoint is written to \code{persisted_state.json} and is sent to the amplifier after a later serial reconnect.

\section{Amplifier warning conditions}
\begin{itemize}
\item A gain warning opens when \code{abs(gain_delta)} is greater than the configured gain
\end{itemize}
tolerance.

\begin{itemize}
\item Minimum and maximum thresholds can be set independently for PiA, PoA, PiB, PoB,
\end{itemize}
and temperature.

\begin{itemize}
\item An empty boundary disables that boundary.
\item For the same field, a minimum must be lower than its maximum.
\end{itemize}
Changing a threshold does not erase historical warning events. It changes future evaluation and can cause the current condition to clear or open on the next measurement.

\chapter{FTS-LS laser station operation}
\section{Polling and command session}
The application assumes that the station CLI accepts the implemented \code{show} commands and returns colon-separated fields. This must be verified with the real station. The worker normally polls every \code{FTS_LS_POLL_SECONDS} (default 10 seconds) and displays the values returned by the CLI.

Operator commands enter a queue with a maximum of 32 items. A command timeout or full queue does not prove that the station did not act. Always inspect the next complete status snapshot before retrying a disruptive command.

\section{How to interpret status labels}
The panel displays state text returned by the CLI but does not classify it as good or bad. \code{--} means that no value was received. The accepted state names and their meanings remain an implementation assumption until the CLI response is verified against the laser-station manual and real hardware.

\section{Internal laser}
The live card shows state, optical frequency, wavelength, central-frequency setpoint, and scanning span. Operator controls include:

\begin{itemize}
\item laser power on/off;
\item central frequency inside the configured safe GHz range;
\item \code{normal} or \code{central-frequency} mode;
\item frequency span from 100 to 10000 MHz;
\item force re-lock.
\end{itemize}
Power, central frequency, mode, and force re-lock require confirmation because they may interrupt transfer. A successful HTTP response means the console command completed; the next poll confirms the resulting station state.

\section{TEC heater}
The card shows TEC state, set/read temperature, and power usage. Operators can turn TEC power on/off and select a 0-100 deg C setpoint. After a change, monitor the read temperature and station state until stable.

\section{10 MHz synthesizer and reference}
The panel shows synthesizer state, reference source, and external 10 MHz status. The external-reference control permits or forbids use of the external frequency. State text is displayed without additional interpretation.

\section{UL and P1-P7 modules}
All seven physical positions remain visible. Empty positions are not removed or renumbered, preventing P4 data from being mistaken for P3 data.

Each equipped card can show:

\begin{itemize}
\item type and state;
\item optical input or estimated optical-power indication;
\item LF/HF noise;
\item jitter;
\item equivalent distance;
\item description;
\item physical connector identifiers.
\end{itemize}
Available settings depend on module type:

\begin{longtable}{|>{\centering\arraybackslash}m{0.20\linewidth}|>{\centering\arraybackslash}m{0.35\linewidth}|>{\centering\arraybackslash}m{0.35\linewidth}|}
\hline
\thead{Setting} & \thead{Applicable target} & \thead{Allowed values} \\
\hline
\endhead
Description & UL and equipped P1-P7 & Up to 120 characters, no newline \\
\hline
Optical power & P1-P7 & On/off; confirmation required \\
\hline
Downlink distance & Downlink modules & 10-2000 km \\
\hline
Additional NC gain & Downlink modules & 0, 12, or 24 dB \\
\hline
Polarization control & UL or feedback-capable module & On/off \\
\hline
Polarization speed & UL or feedback-capable module & Fast/slow \\
\hline
Polarization mode & UL or feedback-capable module & Continuous/triggered \\
\hline
\end{longtable}
A dirty settings section indicates an edited value that has not yet been saved. Save transmits the changed commands; Cancel restores the last complete snapshot.

\section{FTS-LS alarms}
The current version does not generate equipment alarms from FTS-LS state text or measurements. No undocumented thresholds are applied. Alarm rules may be added after their conditions and limits are confirmed in the laser-station manual or on real hardware. The Warnings screen therefore applies only to the amplifier profile.

\section{Administrative station commands}
\begin{longtable}{|>{\centering\arraybackslash}m{0.28\linewidth}|>{\centering\arraybackslash}m{0.64\linewidth}|}
\hline
\thead{Command} & \thead{Effect and required caution} \\
\hline
\endhead
Reboot management system & Restarts station software; expect the serial session to drop and reconnect \\
\hline
Reset complete station & Resets station power and interrupts transfer; use a maintenance window \\
\hline
Restore optical defaults & Replaces current optical settings; record configuration before use \\
\hline
\end{longtable}
Only Administrators can execute these commands. Do not repeat a reset while the previous operation is unresolved.

\chapter{Warnings and alarm lifecycle}
This chapter applies to the amplifier profile. The current FTS-LS profile does not generate equipment alarms.

\section{Active warning states}
\begin{longtable}{|>{\centering\arraybackslash}m{0.28\linewidth}|>{\centering\arraybackslash}m{0.64\linewidth}|}
\hline
\thead{State} & \thead{Meaning} \\
\hline
\endhead
\code{Open} & The triggering condition is currently present and has not been acknowledged \\
\hline
\code{Acknowledged} & An Operator or Administrator confirmed seeing the warning; the condition remains present \\
\hline
No active warnings & No currently evaluated condition exceeds its rule \\
\hline
\end{longtable}
Acknowledgement does not clear a warning, change a threshold, or repair the device. The warning clears only when a later complete measurement no longer matches the condition.

\section{Historical event states}
\begin{longtable}{|>{\centering\arraybackslash}m{0.28\linewidth}|>{\centering\arraybackslash}m{0.64\linewidth}|}
\hline
\thead{State} & \thead{Meaning} \\
\hline
\endhead
\code{Opened} & A condition changed from absent to present \\
\hline
\code{Cleared} & A previously open condition disappeared \\
\hline
\end{longtable}
The cleared event retains the original opening time and, when available, the duration. Repeated samples while a warning remains active update its current value but do not create a new opening event.

\section{Warning data columns}
\begin{itemize}
\item \textbf{Time} - opening or event timestamp.
\item \textbf{Field} - canonical field, module, or supply.
\item \textbf{Value} - measured/current value.
\item \textbf{Target} - configured threshold or expected state.
\item \textbf{Delta} - difference from target when numeric.
\item \textbf{State} - Open, Acknowledged, Opened, or Cleared.
\item \textbf{Message} - readable cause.
\end{itemize}
\section{Warning-history availability states}
\begin{longtable}{|>{\centering\arraybackslash}m{0.28\linewidth}|>{\centering\arraybackslash}m{0.64\linewidth}|}
\hline
\thead{Message} & \thead{Meaning} \\
\hline
\endhead
\code{No warning events in this range} & Log was read successfully but no event matched filters \\
\hline
\code{System warning log is not available} & The configured syslog export file could not be opened \\
\hline
Rotated-log warning & At least one rotated file could not be read; results may be incomplete \\
\hline
\end{longtable}
Active warnings live in process memory. Historical events are read from the syslog file and its rotations. After an application restart, active warnings are rebuilt only from new device data.

\chapter{History, statistics, and export}
Both profiles use the same history controls, local SQLite database, retention setting, database diagnostics, charts, statistics view, and CSV export. Only the fields recorded for each device are different.

\section{Amplifier history}
The selectable ranges are 5 minutes, 1 hour, 24 hours, 7 days, 30 days, all data, and custom ranges where offered. Charts can be downsampled to \code{HISTORY_MAX_POINTS}, so a chart is a visual representation rather than a promise that every raw row is drawn.

Statistics use raw boundary data and stored hourly aggregates where appropriate. The source line reports the number of raw samples used. \code{Could not load statistics} means the statistics request failed; it does not automatically mean that acquisition stopped.

CSV export streams raw amplifier records from an independent read-only SQLite connection. The delimiter is a semicolon and the first line is \code{sep=;} for common spreadsheet compatibility.

\section{FTS-LS history}
The same time ranges and history actions are available for FTS-LS. The history contains the station and module values received during successful polls. Charts and statistics use numeric values; unavailable or non-numeric fields are omitted. CSV files use the same semicolon delimiter and \code{sep=;} header as amplifier exports. Because no FTS-LS alarm limits are currently defined, the statistics view does not classify values as inside or outside limits.


\chapter{Administration screens}
\section{Access control}
User status is either \textbf{Active} or inactive (checkbox cleared). An inactive user cannot log in even if RADIUS accepts the password. Roles control panel actions, not RADIUS permissions. The local list must retain at least one usable Administrator according to site procedure.

The Syslog Download action exports the current audit/warning log. Redact secrets, session data, community strings, and device credentials before sharing logs.

\section{SNMP configuration}
\begin{longtable}{|>{\centering\arraybackslash}m{0.28\linewidth}|>{\centering\arraybackslash}m{0.64\linewidth}|}
\hline
\thead{State} & \thead{Meaning} \\
\hline
\endhead
Agent enabled & The control panel attempts to run the SNMP responder on the configured fixed port \\
\hline
Agent disabled & No control-panel SNMP response is expected \\
\hline
Waiting for SNMP data & No supported live fields are currently available to display \\
\hline
Could not load SNMP data & Browser request failed; check session and service logs \\
\hline
\end{longtable}
The default SNMP UDP port is 1161. The community must contain at least 12 characters. Trap host/port specify the receiver for warning-open traps. Enabling SNMP does not confirm firewall reachability or trap delivery.

\section{Network configuration}
The current connection section shows NetworkManager data:

\begin{longtable}{|>{\centering\arraybackslash}m{0.28\linewidth}|>{\centering\arraybackslash}m{0.64\linewidth}|}
\hline
\thead{Field/state} & \thead{Meaning} \\
\hline
\endhead
Status & Link/device state reported by NetworkManager, commonly connected, disconnected, unavailable, or unknown \\
\hline
Mode \code{dhcp} & Address is obtained automatically \\
\hline
Mode \code{static} & Address, netmask, gateway, and optional DNS are configured manually \\
\hline
\code{None} & No value is configured/reported for that field \\
\hline
Backend unavailable/not supported & Host agent or \code{nmcli} is unavailable; saving is disabled \\
\hline
\end{longtable}
Only the interface carrying the current administrator session can be changed. Use the stable \code{.local} address before an IP-changing operation. The host creates a NetworkManager checkpoint and requires browser confirmation. If confirmation does not arrive within approximately 60 seconds, the host restores the previous settings automatically.

Messages progress through Loading, Applying, Confirming, Saved, or an error that states the old configuration will be restored. Loss of the page during the checkpoint is not proof of failure; wait for rollback before making another change.

\section{Time diagnostics}
\begin{longtable}{|>{\centering\arraybackslash}m{0.28\linewidth}|>{\centering\arraybackslash}m{0.64\linewidth}|}
\hline
\thead{Field/state} & \thead{Meaning} \\
\hline
\endhead
Reachable \code{Yes} & A direct UDP NTP query returned a valid response \\
\hline
Reachable \code{No} & Primary and fallback query attempts failed \\
\hline
Stratum 0 & Unsynchronized/kiss-o'-death response; not a usable time source \\
\hline
Stratum 1 & Primary reference clock \\
\hline
Stratum 2 or greater & Secondary reference at the displayed level \\
\hline
Leap: \code{No warning} & Server does not announce a leap alarm \\
\hline
Leap: 61/59 seconds & Server announces a leap-second condition \\
\hline
Leap: \code{Not synchronized (alarm)} & Server sets leap indicator 3; do not treat it as synchronized \\
\hline
Clock offset & Estimated server minus local clock offset for this query \\
\hline
Round-trip delay & Measured query/response network time \\
\hline
Root delay/dispersion & NTP server quality indicators \\
\hline
\end{longtable}
The message \code{Synchronized with <server>} means the panel received and decoded a direct NTP response. It is not independent proof that \code{systemd-timesyncd} has disciplined the host clock. Check the leap indicator, stratum, offset, and host service logs when synchronization quality matters.

\section{Service diagnostics - data acquisition}
\begin{longtable}{|>{\centering\arraybackslash}m{0.28\linewidth}|>{\centering\arraybackslash}m{0.64\linewidth}|}
\hline
\thead{State} & \thead{Meaning} \\
\hline
\endhead
\code{CONNECTED} & Serial session open \\
\hline
\code{DISCONNECTED} & Serial session closed or reconnecting \\
\hline
Last error \code{None} & No current serial exception recorded \\
\hline
\code{Switching to /dev/...} & Administrator selected another serial port; reconnect is in progress \\
\hline
\end{longtable}
Only currently discovered \code{/dev/ttyACM*} and \code{/dev/ttyUSB*} devices can be selected from the web settings page. A persistent by-id path can be configured by the package configurator when appropriate.

\section{Service diagnostics - history storage}
\begin{longtable}{|>{\centering\arraybackslash}m{0.28\linewidth}|>{\centering\arraybackslash}m{0.64\linewidth}|}
\hline
\thead{State} & \thead{Meaning} \\
\hline
\endhead
\code{READY} & SQLite connection is open and history operations can proceed \\
\hline
\code{ERROR} & SQLite could not be initialized; record count is reported as zero \\
\hline
\code{API ERROR} & Browser could not read service diagnostics \\
\hline
Retention \code{UNLIMITED} & No record-count pruning; free disk space remains the real limit \\
\hline
Time \code{Not applicable} & Record limit is disabled \\
\hline
Estimate \code{--} & Insufficient rate/disk data for an estimate \\
\hline
\end{longtable}
\code{Last error} records the latest SQLite/storage operation failure. A non-empty error can coexist briefly with a ready connection if a later nonfatal operation failed. Diagnose the error instead of relying on the color alone.

The size includes \code{measurements.db}, \code{measurements.db-wal}, and \code{measurements.db-shm} when present. The write rate is calculated from recent records.

When a finite record limit is reached, the oldest applicable records are removed.

\section{Service diagnostics - system logging}
\begin{longtable}{|>{\centering\arraybackslash}m{0.28\linewidth}|>{\centering\arraybackslash}m{0.64\linewidth}|}
\hline
\thead{State} & \thead{Meaning} \\
\hline
\endhead
Local \code{ENABLED} & Application sends events to configured local syslog destination \\
\hline
Local \code{DISABLED} & No application syslog messages are intentionally sent \\
\hline
Remote \code{ENABLED} & Rsyslog forwarding is configured \\
\hline
Remote \code{DISABLED} / Not configured & No remote receiver is configured \\
\hline
Heartbeat interval & Periodic online event interval \\
\hline
Heartbeat \code{DISABLED} & Interval equals zero \\
\hline
\end{longtable}
Remote \code{ENABLED} means configured, not delivered. Rsyslog may queue messages while the receiver is offline.

\chapter{Files, storage, retention, and backup}
\section{Installed file locations}
\begin{itemize}
\item \code{/usr/lib/amp-panel/}\\
Installed Python application, static files, templates, and vendored dependencies.

\item \code{/usr/bin/amp-panel}\\
Administration command-line interface.

\item \code{/etc/amp-panel/amp-panel.env}\\
Root-owned configuration and secrets, mode 0600. Preserved during a normal package upgrade.

\item \code{/var/lib/amp-panel/measurements.db}\\
SQLite history for the active device profile. Retained when the package is removed or purged.

\item \code{/var/lib/amp-panel/measurements.db-wal} and
\code{/var/lib/amp-panel/measurements.db-shm}\\
Runtime companions of the SQLite database. Back them up consistently with the main database.

\item \code{/var/lib/amp-panel/persisted_state.json}\\
Panel settings, users, SNMP configuration, and profile-specific settings. Mode 0600 and retained when the package is removed or purged.

\item \code{/var/log/amp-panel/amp-panel.log} and
\code{/var/log/amp-panel/amp-panel.log.*}\\
Current and rotated warning, lifecycle, heartbeat, and audit logs. Rotation controls their retention.

\item \code{/run/amp-panel/network-agent.sock}\\
Restricted runtime socket between the web service and network agent. Recreated after boot or service restart.

\item \code{/etc/rsyslog.d/30-amp-panel.conf} and
\code{/etc/logrotate.d/amp-panel}\\
System logging route and log-rotation policy.

\item \code{/etc/avahi/services/amp-panel.service}\\
Advertisement of the panel's \code{.local} service.

\item \code{/etc/systemd/timesyncd.conf.d/amp-panel.conf}\\
Time-server configuration.

\item \code{/etc/systemd/system/amp-panel.service.d/paths.conf}\\
Generated data and log path permissions for the service sandbox.
\end{itemize}
If \code{AMP_PANEL_DATA_DIR} is changed, the database and state paths move below the selected directory. Allowed locations are \code{/var/lib/amp-panel} or paths below \code{/mnt}, \code{/media}, and \code{/srv}.

\section{What \code{persisted_state.json} contains}
\begin{itemize}
\item last accepted amplifier gain setpoint;
\item amplifier warning limits and gain tolerance;
\item local usernames, roles, and active flags;
\item SNMP enablement, fixed port, community, and trap destination;
\item syslog heartbeat interval, database record limit, and selected serial port.
\end{itemize}
It does not contain RADIUS passwords. It can contain the SNMP community and must be protected as sensitive data.

\section{What SQLite contains}
Both profiles use the same SQLite database file. The internal tables reflect the different data returned by each device:
\begin{itemize}
\item \code{samples} - columnar amplifier samples;
\item \code{setpoint_events} - amplifier gain-set changes;
\item \code{device_snapshots} - normalized FTS-LS snapshots;
\item \code{hourly_statistics} - precomputed amplifier summaries;
\item metadata/count tables and triggers used for efficient diagnostics.
\end{itemize}
Warning history is not stored in SQLite. It comes from syslog. Current active warnings exist only in process memory.

\section{Retention behavior}
\code{DATABASE_MAX_RECORDS=0} means no count-based pruning. It does not reserve disk space and does not prevent a full filesystem. With a positive limit, oldest amplifier samples or FTS-LS snapshots are pruned when the limit is applied. Setpoint events are not included in the amplifier sample limit.

Log retention is independent and follows logrotate configuration. Changing the database record limit does not change warning/audit-log retention.

\section{Safe offline backup}
\begin{lstlisting}
sudo amp-panel stop
sudo sqlite3 /var/lib/amp-panel/measurements.db \
  'PRAGMA wal_checkpoint(TRUNCATE); PRAGMA integrity_check;'
sudo cp -a /etc/amp-panel/amp-panel.env /secure/destination/
sudo cp -a /var/lib/amp-panel/measurements.db /secure/destination/
sudo cp -a /var/lib/amp-panel/persisted_state.json /secure/destination/
sudo cp -a /var/log/amp-panel/amp-panel.log* /secure/destination/   # optional
sudo amp-panel start
sudo amp-panel doctor
\end{lstlisting}
The integrity result must be \code{ok}. If \code{persisted_state.json} does not yet exist, no persistent panel setting has been saved. Store backups containing secrets in encrypted, access-controlled storage.

Never restore a WAL file from a different moment than its main database. For an online copy, treat DB/WAL/SHM as one consistent set or use the SQLite backup API.

\section{Moving the data directory}
\begin{lstlisting}
sudo amp-panel data-dir migrate /mnt/amp-panel-data
\end{lstlisting}
This stops the application, checkpoints and verifies SQLite, copies managed data through temporary files, verifies the destination, updates configuration, and restarts services. The source is intentionally not deleted.

\code{amp-panel data-dir use PATH} only switches to an already prepared directory; it does not copy data.

\chapter{Installation and configuration}
\section{Building the Debian package}
Building is a separate operation from installing the finished package. On a clean
Debian build system, install the required tools once:

\begin{lstlisting}
sudo apt update
sudo apt install build-essential debhelper git python3 python3-pip
\end{lstlisting}

When already logged in as \code{root}, omit \code{sudo} from these commands.

The \code{debhelper} package provides the required
\code{debhelper-compat (= 13)} build dependency. The \code{python3-pip} package is
needed when \code{packaging/vendor} must be prepared from \code{requirements.txt}.

Run the build from the project root:

\begin{lstlisting}
./packaging/build_deb.sh
\end{lstlisting}

The script checks the declared Debian build dependencies before starting. If a
dependency is missing, it stops and prints the installation command. A successful
build writes \code{amp-panel_<version>_<architecture>.deb} to the parent directory
of the project. Build on Linux for the same architecture as the target device.

\section{Package installation}
\begin{lstlisting}
sudo apt install ./amp-panel_<version>_<architecture>.deb
sudo amp-panel configure
\end{lstlisting}
The configurator asks for device profile, web port, serial port, data directory, initial administrator, RADIUS, mDNS hostname, and profile-specific credentials or ranges. It writes \code{/etc/amp-panel/amp-panel.env} atomically with mode 0600 and starts the services.

After installation:

\begin{lstlisting}
sudo amp-panel doctor
sudo amp-panel status
sudo amp-panel paths
sudo amp-panel logs -n 100
\end{lstlisting}
\section{Important configuration groups}
\begin{longtable}{|>{\centering\arraybackslash}m{0.28\linewidth}|>{\centering\arraybackslash}m{0.64\linewidth}|}
\hline
\thead{Group} & \thead{Important variables} \\
\hline
\endhead
Web/device & \code{AMP_PANEL_PORT}, \code{DEVICE_PROFILE}, \code{DEVICE_NAME}, \code{MDNS_HOSTNAME} \\
\hline
Serial & \code{SERIAL_PORT}, \code{SERIAL_BAUDRATE} \\
\hline
Amplifier & \code{GAIN_SET_MIN}, \code{GAIN_SET_MAX} \\
\hline
FTS-LS & \code{FTS_LS_USERNAME}, \code{FTS_LS_PASSWORD}, poll interval, safe frequency range \\
\hline
Storage & \code{AMP_PANEL_DATA_DIR}, \code{DATABASE_FILE}, \code{PERSISTED_STATE_FILE}, record limit \\
\hline
Login & Initial administrator, RADIUS server/port/secret, attempts, session age \\
\hline
Logging & Local/remote syslog, facility, timezone, heartbeat, export file \\
\hline
Integrations & SNMP port/community, NTP primary/fallback, network-agent socket \\
\hline
\end{longtable}
Do not hand-edit paths independently. The database and persistent-state files must remain inside the selected data directory. Prefer \code{sudo amp-panel configure}.

\section{CLI reference}
\begin{longtable}{|>{\centering\arraybackslash}m{0.28\linewidth}|>{\centering\arraybackslash}m{0.64\linewidth}|}
\hline
\thead{Command} & \thead{Purpose} \\
\hline
\endhead
\code{amp-panel configure} & Validate and write host configuration \\
\hline
\code{amp-panel doctor} & Check configuration, paths, SQLite, serial presence, and services \\
\hline
\code{amp-panel status} & Show systemd status without a pager \\
\hline
\code{amp-panel logs -n N} & Show recent application journal lines \\
\hline
\code{amp-panel logs -f} & Follow application journal \\
\hline
\code{amp-panel start/stop/restart} & Control application service \\
\hline
\code{amp-panel paths} & Print configuration, data, log, and runtime locations \\
\hline
\code{amp-panel data-dir migrate PATH} & Safely copy and switch managed data \\
\hline
\code{amp-panel data-dir use PATH} & Switch to an existing prepared data directory \\
\hline
\code{amp-panel version} & Print installed version \\
\hline
\end{longtable}
\chapter{Routine operation and maintenance}
\section{Start-of-shift check}
\begin{enumerate}
\item Sign in and confirm the expected role.
\item Confirm \code{CONNECTED}.
\item Compare Last update with System time.
\item Review active warnings before changing any setting.
\item Check critical amplifier values or FTS-LS lock states.
\item If responsible for administration, confirm history storage is \code{READY} and
\end{enumerate}
inspect free disk space.

\section{Before a control change}
\begin{itemize}
\item Record the starting value and relevant device state.
\item Change one parameter at a time.
\item Confirm the requested range and unit.
\item For FTS-LS, allow at least one full poll after the command.
\item Never interpret a browser notification alone as physical confirmation.
\item Use a maintenance window for power, laser, network, reboot, or default actions.
\end{itemize}
\section{End-of-shift or incident record}
Record incident time, device profile, visible state, last update, active warnings, operator action, resulting state, and any exported log/history. Do not include passwords, RADIUS secrets, SNMP communities, or cookies.

\section{Package upgrade}
\begin{enumerate}
\item Create and verify a backup.
\item Install the new package.
\item Run \code{sudo amp-panel configure} if configuration requirements changed.
\item Run \code{sudo amp-panel doctor}.
\item Verify login, profile, serial connection, new history records, and integrations.
\end{enumerate}
The project supports only the current package configuration and supported database schema upgrades.

\chapter{Troubleshooting by visible state}
\section{Dashboard does not open}
\begin{lstlisting}
sudo amp-panel status
sudo amp-panel logs -n 200
sudo amp-panel paths
\end{lstlisting}
Verify the configured TCP port (1024-65535), host IP, firewall, and Avahi if using \code{.local}. An import-time configuration error prevents service startup.

\section{\code{DISCONNECTED}}
\begin{enumerate}
\item Check device power and cable.
\item Inspect \code{/dev/serial/by-id}, \code{/dev/ttyACM*}, and \code{/dev/ttyUSB*}.
\item Confirm the configured port and baud rate.
\item Confirm the \code{amp-panel} user has serial-device group access.
\item Check whether another program owns the port.
\item For FTS-LS, verify console username/password and expected prompt.
\item Follow logs while reconnecting.
\end{enumerate}
\begin{lstlisting}
sudo amp-panel doctor
sudo amp-panel logs -f
\end{lstlisting}
\section{\code{CONNECTED} but values are \code{--}}
The port is open but no supported complete message has been accepted. Check Last update, raw device output, line termination, firmware labels, FTS-LS prompt, and protocol parsing. Do not change frontend labels to compensate for a raw protocol change; update the adapter mapping.

\section{\code{API ERROR}}
Check browser connectivity and session, then \code{amp-panel status} and journal logs. If only one page fails, inspect that subsystem: database, network agent, NTP, or syslog. A serial device can continue operating while a browser request fails.

\section{History storage \code{ERROR}}
\begin{enumerate}
\item Stop the service.
\item Back up every \code{measurements.db*} file.
\item Check disk free space, permissions, and storage-device errors.
\item Run SQLite integrity checking on a copy or with the service stopped.
\item Do not reduce retention or migrate the directory until integrity is understood.
\end{enumerate}
\section{Active warnings disappear after restart}
This is expected because active warnings are in memory. New device measurements rebuild the list. Historical Opened/Cleared events remain in syslog rotations if logging and retention are working.

\section{Warning history unavailable}
Check \code{/var/log/amp-panel/amp-panel.log}, rotations, \code{SYSLOG_EXPORT_FILE}, rsyslog, log permissions, and membership of the service in group \code{adm}.

\section{NTP reachable \code{No}}
Verify UDP 123, DNS, primary/fallback addresses, timeout, and server status. If Reachable is Yes but leap indicator says Not synchronized, choose a healthy time source. A small offset is not sufficient proof when stratum/leap status is bad.

\section{Network change rolled back}
Rollback is the safe result when the new address could not be confirmed from the same active interface. Reconnect to the old address, check the network-agent service, use the \code{.local} URL, and retry only after confirming the original state.

\section{SNMP unavailable}
Confirm SNMP is enabled, UDP port 1161 (or configured fixed port), community, firewall, and current live data. Test from another host if possible. Trap configuration and polling configuration are separate.

\section{Information required for a problem report}
\begin{itemize}
\item output of \code{amp-panel version} and \code{amp-panel doctor};
\item device profile and selected serial port;
\item incident time and timezone;
\item visible status and Last update;
\item last 200 relevant journal lines;
\item recent change that preceded the issue;
\item sanitized log/history samples.
\end{itemize}
\chapter{Technical reference}
\section{Canonical boundaries}
Hardware-specific parsing and command spelling belong in:

\begin{itemize}
\item \code{app/protocols/amplifier.py} for the amplifier;
\item \code{app/protocols/fts_ls.py} for the laser station.
\end{itemize}
Canonical field definitions live in \code{app/core/device_schema.py}. The API, database, warning logic, SNMP, and frontend should consume canonical names. A firmware rename should normally require an adapter change plus its protocol test, not changes throughout the application.

\section{Main application areas}
\begin{longtable}{|>{\centering\arraybackslash}m{0.28\linewidth}|>{\centering\arraybackslash}m{0.64\linewidth}|}
\hline
\thead{Area} & \thead{Location} \\
\hline
\endhead
Configuration and state & \code{app/core/} \\
\hline
Device protocol adapters & \code{app/protocols/} \\
\hline
HTTP endpoints & \code{app/api/} \\
\hline
Serial, database, syslog, SNMP, NTP, network bridge & \code{app/services/} \\
\hline
Browser structure & \code{templates/index.html} \\
\hline
Browser behavior & \code{static/js/dashboard*.js} \\
\hline
Package administration CLI & \code{tools/amp_panel_cli.py} \\
\hline
Restricted host network agent & \code{tools/network_agent.py} \\
\hline
Debian/systemd packaging & \code{debian/} and \code{packaging/systemd/} \\
\hline
\end{longtable}
\section{Safety rules}
\begin{enumerate}
\item Treat a fresh physical status poll as confirmation of a device change.
\item Never use factory default without a recorded recommissioning plan.
\item Never copy a live SQLite main file without its WAL/SHM context.
\item Never publish configuration, persisted state, or logs without redacting
\end{enumerate}
credentials and communities.

\begin{enumerate}
\item Never bypass the network checkpoint/confirmation mechanism.
\item Preserve canonical application names when adapting firmware differences.
\item Validate all final ranges against production hardware documentation before
\end{enumerate}
commissioning.

\section{Quick state reference}
\begin{longtable}{|>{\centering\arraybackslash}m{0.28\linewidth}|>{\centering\arraybackslash}m{0.64\linewidth}|}
\hline
\thead{Visible state} & \thead{Immediate interpretation} \\
\hline
\endhead
\code{CONNECTED} & Serial session open \\
\hline
\code{DISCONNECTED} & Serial session closed/reconnecting \\
\hline
\code{API ERROR} & Browser request to the control panel failed \\
\hline
\code{READY} & SQLite available \\
\hline
\code{ERROR} & SQLite initialization failed \\
\hline
Warning \code{Open} & Condition active and unacknowledged \\
\hline
Warning \code{Acknowledged} & Condition still active but seen \\
\hline
History \code{Opened} & Warning began \\
\hline
History \code{Cleared} & Warning condition ended \\
\hline
Syslog \code{ENABLED} & Logging/forwarding configured, not delivery proof \\
\hline
Retention \code{UNLIMITED} & No count limit; disk capacity still applies \\
\hline
\end{longtable}
\textbf{End of manual}

\backmatter
\end{document}
